\documentclass[UTF8]{ctexart}
\usepackage{xeCJK}
\usepackage{geometry}
\geometry{left=2cm,right=2cm,top=2.5cm,bottom=2.5cm}
\setlength{\headheight}{12.7pt}

\usepackage{titlesec}
\titleformat{\section}
  {\normalfont\large\bfseries}
  {\thesection}
  {1em}
  {}
\titlespacing*{\section}
  {0pt}
  {3.5ex plus 1ex minus .2ex}
  {2.3ex plus .2ex}

\usepackage{amsmath}
\usepackage{amssymb}
\usepackage{amsthm}
\usepackage{enumitem}
\setlist[enumerate]{itemsep=0pt,topsep=0pt,parsep=0pt,leftmargin=0pt,itemindent=2.5em}
\setlist[itemize]{itemsep=0pt,topsep=0pt,parsep=0pt,leftmargin=0pt,itemindent=2.5em}
\usepackage{fancyhdr}
\pagestyle{fancy}
\fancyhf{}
\usepackage{graphicx}
\usepackage{hyperref}
\usepackage{indentfirst}
\usepackage{lmodern}


\begin{document}
\setlength{\abovedisplayskip}{1pt}
\setlength{\belowdisplayskip}{1pt}

\section{拓扑基}

\hypertarget{sec:定义1.1}{}
\noindent\textbf{定义1.1 (拓扑基)}

给定集合$X$和它的子集族$\mathcal{B}\subset\mathcal{P}(X)$. 如果:
\vspace{3pt}
\begin{enumerate}
    \item $\displaystyle\bigcup\mathcal{B}=X$,
    \vspace{3pt}
    \item 对任意的$B_1,B_2\in\mathcal{B}$和$x\in B_1\cap B_2$,
    存在$B_3\in\mathcal{B}$使得$x\in B_3\subset B_1\cap B_2$,
\end{enumerate}
就称$\mathcal{B}$是$X$上的一个\textbf{拓扑基}.

\vspace{10pt}

\hypertarget{sec:定义1.2}{}
\noindent\textbf{定义1.2 (生成拓扑)}

任给集合$X$和$X$上的拓扑基$\mathcal{B}$,
\[
    \langle\mathcal{B}\rangle=\{U\subset X\mid
    \forall x\in U\,\,\exists B\in\mathcal{B}:\,x\in B\subset U\}
\]
是$X$上包含$\mathcal{B}$的最粗拓扑, 称作由$\mathcal{B}$生成的拓扑.
也称$\mathcal{B}$是$\langle\mathcal{B}\rangle$的一个基.

\noindent{\kaishu 验证.}

显然$\langle\mathcal{B}\rangle$包含$\mathcal{B}$.
首先证明$\langle\mathcal{B}\rangle$是拓扑.

\begin{enumerate}
    \item 显然$\varnothing\in\langle\mathcal{B}\rangle$.
    因为$\displaystyle\bigcup\mathcal{B}=X$, 所以$X\in\langle\mathcal{B}\rangle$.
    \vspace{3pt}
    \item 对$\langle\mathcal{B}\rangle$中任意的$(U_i)_{i\in I}$,
    任取$\displaystyle x\in\bigcup_{i\in I}U_i$, 存在$i\in I$使得$x\in U_i$,
    依定义存在$B\in\mathcal{B}$使得
    \vspace{3pt}
    \[
        x\in B\subset U_i\quad\implies\quad
        x\in B\subset U_i\subset\bigcup_{i\in I}U_i,\\[-3pt]
    \]
    所以$\displaystyle\bigcup_{i\in I}U_i\in\langle\mathcal{B}\rangle$.
    \vspace{3pt}
    \item 对任意的$U_1,U_2\in\langle\mathcal{B}\rangle$, 任取$x\in U_1\cap U_2$,
    依定义存在$B_1,B_2\in\mathcal{B}$使得
    \[
        x\in B_1\subset U_1\quad\land\quad x\in B_2\subset U_2,
    \]
    此时存在$B_3\in\mathcal{B}$使得
    \[
        x\in B_3\subset B_1\cap B_2\subset U_1\cap U_2,
    \]
    所以$U_1\cap U_2\in\langle\mathcal{B}\rangle$.
\end{enumerate}

\vspace{3pt}
接下来, 如果拓扑$\mathcal{T}$包含$\mathcal{B}$,
那么对任意的$U\in\langle\mathcal{B}\rangle$,
对任意的$x\in U$选取一个$B_x\in\mathcal{B}$使得
\[
    x\in B_x\subset U,
\]
易证$\displaystyle U=\bigcup_{x\in U}B_x$. 因为每个$B_x\in\mathcal{T}$,
所以$U\in\mathcal{T}$. 所以$\langle\mathcal{B}\rangle\subset\mathcal{T}$,
即$\langle\mathcal{B}\rangle$是$X$上包含$\mathcal{B}$的最粗拓扑. \hfill $\qedsymbol$

\vspace{10pt}

\hypertarget{sec:定理1.1}{}
\noindent\textbf{定理1.1}

给定拓扑空间$(X,\mathcal{T})$和$\mathcal{B}\subset\mathcal{T}$, 以下等价:
\begin{enumerate}
    \item $\mathcal{B}$是此空间的基.
    \item 对任意开集$U$和点$x\in U$, 存在$B\in\mathcal{B}$使得$x\in B\subset U$.
    \vspace{3pt}
    \item 对任意开集$U$, 存在$\mathcal{B}$中的$(B_i)_{i\in I}$使得
    $\displaystyle U=\bigcup_{i\in I}B_i$.
\end{enumerate}

% 毕业论文终于大概定稿了. 久违地写一点数学.
% 儿童节快乐.
% 2026.06.01





\newpage

\section{拓扑子基}

\hypertarget{sec:定义2.1}{}
\noindent\textbf{定义2.1 (拓扑子基)}

给定集合$X$和$X$的子集族$\mathcal{S}\subset\mathcal{P}(X)$.
如果$\displaystyle\bigcup\mathcal{S}=X$,
就称$\mathcal{S}$是$X$上的一个\textbf{拓扑子基}.

\vspace{10pt}

拓扑子基通过有限交扩张操作可以得到拓扑基, 从而可以生成拓扑.

\vspace{10pt}

\hypertarget{sec:定义2.2}{}
\noindent\textbf{定义2.2 (生成滤子)}

任给集合$X$和$X$上的拓扑子基$\mathcal{S}$,
\vspace{3pt}
\[
    \mathcal{B}=\left\{\,\bigcap\mathcal{R}\,\,\middle\vert\,\,
    \mathcal{R}\text{\,是\,}\mathcal{S}\text{\,的非空有限子集\,}\right\}\\[3pt]
\]
是$X$上包含$\mathcal{S}$的拓扑基.

进一步, $\langle\mathcal{B}\rangle$是$X$上包含$\mathcal{S}$的最粗拓扑,
也记作$\langle\mathcal{S}\rangle$, 称作由$\mathcal{S}$生成的拓扑,
称$\mathcal{S}$是其子基.

\noindent{\kaishu 验证.}

显然$\mathcal{B}$包含$\mathcal{S}$. 首先证明$\mathcal{B}$是拓扑基.

% 又到了高考的时候, 祝所有人高考顺利.
% 我现在可能与他们同样地紧张. 也希望自己三天后的答辩顺利.
% 2026.06.07

\begin{enumerate}
    \item $\displaystyle\bigcup\mathcal{B}\supset\bigcup\mathcal{S}=X$,
    从而$\displaystyle\bigcup\mathcal{B}=X$.
    \vspace{3pt}
    \item 对任意的$B_1,B_2\in\mathcal{B}$, 存在$\mathcal{S}$的非空有限子集
    $\mathcal{R}_1,\mathcal{R}_2$使得
    \[
        B_1=\bigcap\mathcal{R}_1\quad\land\quad B_2=\bigcap\mathcal{R}_2,
    \]
    从而$\displaystyle B_1\cap B_2=\bigcap\,(\mathcal{R}_1\cup\mathcal{R}_2)$
    也属于$\mathcal{B}$.
\end{enumerate}

\vspace{3pt}
接下来, 如果拓扑$\mathcal{T}$包含$\mathcal{S}$,
那么对任意的$B\in\mathcal{B}$, 存在$\mathcal{S}$的非空有限子集$\mathcal{R}$使得
\[
    B=\bigcap\mathcal{R}\in\mathcal{T},
\]
所以$\mathcal{B}\subset\mathcal{T}$,
从而$\langle\mathcal{B}\rangle\subset\mathcal{T}$,
即$\langle\mathcal{B}\rangle$是$X$上包含$\mathcal{S}$的最粗拓扑. \hfill $\qedsymbol$

% 答辩结束了, 不知道结果如何.
% 现在也并没有之前想的那样放松. 有点慌...
% 2026.06.10

\vspace{10pt}

\hypertarget{sec:定理2.1}{}
\noindent\textbf{定理2.1}

给定拓扑空间$(X,\mathcal{T})$和$\mathcal{S}\subset\mathcal{T}$, 以下等价:
\begin{enumerate}
    \item $\mathcal{S}$是此空间的子基.
    \vspace{3pt}
    \item 对任意开集$U$, 存在$\mathcal{S}$的一族非空有限子集$(\mathcal{R}_i)_{i\in I}$
    使得$\displaystyle U=\bigcup_{i\in I}\bigcap\mathcal{R}_i$.
\end{enumerate}

\end{document}